\begin{resumo}
A Hipótese de Mercados Eficientes tem sido revisitada pelas Finanças Comportamentais, que destacam o papel do fluxo informacional e da racionalidade limitada dos investidores na precificação de ativos. Em mercados emergentes, ativos de alta liquidez e forte exposição política, como as ações preferenciais da Petrobras (PETR4), exibem sensibilidade pronunciada a choques de notícias. Esta dissertação investiga o impacto do \textbf{sentimento} de notícias financeiras, definido como o escore de polaridade no intervalo $[-1,+1]$ atribuído a um texto por um classificador, na previsão da direção e da volatilidade do ativo PETR4 na B3, no período de 2018 a 2025. Adota-se uma abordagem empírico-quantitativa estruturada em um \textit{pipeline} reprodutível. As notícias foram coletadas diretamente de cinco portais por meio de sua \textit{Application Programming Interface} (API) REST, o que assegurou a \textbf{data e a hora exatas} de cada publicação, requisito indispensável ao alinhamento causal \textit{Lead-Lag} entre a notícia e o pregão. O corpus resultante, de 205.697 notícias, foi organizado por uma taxonomia supervisionada de sete categorias temáticas e teve seu sentimento extraído pelo \textbf{FinBERT-PT-BR}, modelo da família BERT ajustado ao domínio financeiro em português. A volatilidade foi modelada por um GARCH(1,1) com distribuição \textit{t}-Student, e a previsão de direção empregou \textit{Support Vector Machines} (SVM) e \textit{Extreme Gradient Boosting} (XGBoost) sob uma estrutura de \textit{Data Fusion} com divisão cronológica tripla. Os resultados mostram que a fusão do sentimento com os preços eleva a acurácia de teste de 50,9\% (linha de base) para 54,5\%, ganho estatisticamente superior ao acaso (teste binomial, $p=0{,}024$). A causalidade de Granger confirma que o sentimento antecede o retorno em um dia e, sobretudo, antecede a volatilidade de forma forte e persistente. Conclui-se que o sentimento contribui de modo modesto, porém significativo, para a previsão de direção, e de modo mais nítido para a volatilidade, em consonância com a quase-eficiência do mercado.

\vspace{\onelineskip}
\noindent
\textbf{Palavras-chave}: Finanças Comportamentais. Processamento de Linguagem Natural. FinBERT-PT-BR. Aprendizado de Máquina. GARCH. PETR4.
\end{resumo}
